\documentclass[10pt,a4paper]{amsart}
\usepackage[margin=19mm]{geometry}
\usepackage{amsmath,amssymb,mathtools}
\usepackage[scr=rsfs,cal=euler]{mathalfa}
\usepackage{xfrac}
\usepackage[unicode,hidelinks,bookmarks=false]{hyperref}
\newcommand{\ie}{i.e.\ }
\newcommand{\cf}{cf.\ }
\newcommand{\resp}{respectively\ }
\newcommand{\Subsection}{\subsection}
\providecommand{\nobreakdash}{\nobreak\hbox{-}}
\setlength{\emergencystretch}{2em}
\multlinegap0pt
\usepackage{mathtools}
\newtheorem{theorem}{Theorem}
\newtheorem{lemma}[theorem]{Lemma}
\newtheorem{proposition}[theorem]{Proposition}
\newtheorem{corollary}[theorem]{Corollary}
\newtheorem{definition}[theorem]{Definition}
\newtheorem{question}{Question}
\renewcommand{\thequestion}{}
\newtheorem{remark}[theorem]{Remark}
\DeclareMathOperator{\TR}{TR}
\newcommand{\R}{\mathrm{R}}
\begin{document}
\title[Connected Counterexamples for Target Ramsey Numbers]{Connected Counterexamples for Target Ramsey Numbers}
\author{Guillaume Lecomte}
\date{}
\maketitle
\noindent\emph{Source and licence.} Connected Counterexamples for Target Ramsey Numbers. Version arXiv:2608.06446v1 (CC BY 4.0). This material is distributed under the Creative Commons Attribution 4.0 International licence, http://creativecommons.org/licenses/by/4.0/, as recorded in the arXiv OAI licence metadata for this exact version and shown on the abstract page https://arxiv.org/abs/2608.06446v1. Full rights notice: licenses/arxiv-2608.06446-CC-BY-4.0.txt.\par\smallskip
\noindent\textit{Editorial scope.} Counted unit: the original \texttt{theorem} environment \texttt{thm:dgrid} at source lines 188--202 together with its complete proof at lines 204--210; the delivered document keeps source lines 71--211 so that every referenced label and every citation resolves inside it. Native editable source is the paper's own concatenated TeX; the AMS wrapper is editorial formatting only. No publisher class, upstream script or unrelated artwork is redistributed.
\begin{lemma}\label{lem:count}
For every graph $G$ without isolated vertices,
\[
\TR(G) \geq e(G) + 1.
\]
\end{lemma}

\begin{proof}
Set $m = e(G)$, and suppose that a Ramsey chain with target $G$ occurs in $K_n$. Since its links are pairwise edge-disjoint and have sizes $1, 2, \ldots, m$, their union contains exactly
\[
\sum_{i=1}^{m} i = \binom{m+1}{2}
\]
distinct edges. Hence
\[
\binom{n}{2} \geq \binom{m+1}{2}.
\]
The map $s \mapsto \binom{s}{2}$ is strictly increasing on the positive integers, so $n \geq m+1$. Equivalently, for $n \leq m$ no colouring of $K_n$ contains a Ramsey chain with target $G$, and the conclusion follows from the definition of $\TR(G)$.
\end{proof}

\begin{proposition}\label{prop:criterion}
If a graph $G$ without isolated vertices satisfies
\[
\R(G) \leq e(G),
\]
then $\TR(G) > \R(G)$.
\end{proposition}

\begin{proof}
Lemma~\ref{lem:count} gives $\TR(G) \geq e(G) + 1 > \R(G)$.
\end{proof}

Thus it suffices to find graphs whose ordinary Ramsey number is at most their number of edges. Large grids provide a particularly structured family of this kind.

\section{Connected planar counterexamples}

Let $P_t$ denote the path on $t$ vertices and let
\[
\Gamma_t = P_t \square P_t
\]
be the square $t \times t$ grid. For $t \geq 3$, the graph $\Gamma_t$ is connected, planar, bipartite and has maximum degree four. Moreover,
\begin{equation}\label{eq:grid}
v(\Gamma_t) = t^2, \qquad e(\Gamma_t) = 2t(t-1) = 2t^2 - 2t.
\end{equation}

Mota, S\'ark\"ozy, Schacht and Taraz proved that the two-colour Ramsey number of the grid $G_{a,b} = P_a \square P_b$ satisfies
\[
\R(G_{a,b}) = \left( \tfrac{3}{2} + o(1) \right) ab,
\]
where the $o(1)$ tends to zero as $ab \to \infty$; see \cite[Corollary 1.4]{MSST}. In particular, as $t \to \infty$,
\begin{equation}\label{eq:ramseygrid}
\R(\Gamma_t) = \left( \tfrac{3}{2} + o(1) \right) t^2.
\end{equation}
Only the upper bound implicit in \eqref{eq:ramseygrid} will be used.

\begin{theorem}\label{thm:planar}
For every sufficiently large integer $t$,
\[
\TR(\Gamma_t) > \R(\Gamma_t).
\]
Consequently, there are infinitely many connected planar bipartite graphs of maximum degree four whose target Ramsey number is strictly larger than their ordinary Ramsey number.
\end{theorem}

\begin{proof}
By Lemma~\ref{lem:count} and \eqref{eq:grid},
\begin{equation}\label{eq:trgrid}
\TR(\Gamma_t) \geq 2t^2 - 2t + 1.
\end{equation}
On the other hand, \eqref{eq:ramseygrid} implies, for example, that
\[
\R(\Gamma_t) \leq \tfrac{7}{4} t^2
\]
for all sufficiently large $t$. Since
\[
2t^2 - 2t + 1 - \tfrac{7}{4} t^2 = \tfrac{1}{4} t^2 - 2t + 1 \longrightarrow +\infty,
\]
we have
\[
\R(\Gamma_t) < 2t^2 - 2t + 1 \leq \TR(\Gamma_t)
\]
for every sufficiently large $t$.
\end{proof}

\begin{corollary}
As $t \to \infty$,
\[
\TR(\Gamma_t) - \R(\Gamma_t) \geq \left( \tfrac{1}{2} - o(1) \right) t^2,
\]
and
\[
\liminf_{t \to \infty} \frac{\TR(\Gamma_t)}{\R(\Gamma_t)} \geq \frac{4}{3}.
\]
\end{corollary}

\begin{proof}
Combine \eqref{eq:trgrid} with \eqref{eq:ramseygrid}, and divide by $t^2$ for the ratio statement.
\end{proof}

\begin{remark}[No explicit first grid]
The asymptotic theorem quoted in \eqref{eq:ramseygrid} does not by itself identify the least $t$ for which $\Gamma_t$ is a counterexample. Determining the smallest connected counterexample, or even the first square grid satisfying $\TR(\Gamma_t) > \R(\Gamma_t)$, is a separate finite problem.
\end{remark}

\section{Higher-dimensional grids}

For integers $d \geq 2$ and $t \geq 2$, let
\[
\Gamma_{d,t} = P_t^{\square d}
\]
be the $d$-dimensional grid. It is connected and bipartite, its two vertex classes differ in size by at most one, and
\begin{equation}\label{eq:dgrid}
v(\Gamma_{d,t}) = t^d, \qquad e(\Gamma_{d,t}) = d\, t^{d-1}(t-1), \qquad \Delta(\Gamma_{d,t}) = 2d.
\end{equation}
Its bandwidth is at most $t^{d-1} = o(t^d)$ for fixed $d$. Hence, for fixed $d$, the two-colour result of Mota, S\'ark\"ozy, Schacht and Taraz applies and gives
\begin{equation}\label{eq:ramseydgrid}
\R(\Gamma_{d,t}) \leq \left( \tfrac{3}{2} + o(1) \right) t^d \qquad (t \to \infty,\ d \text{ fixed});
\end{equation}
see \cite[Theorem 1.2 and the discussion of higher-dimensional grids]{MSST}.


\begin{theorem}\label{thm:dgrid}
For every fixed integer $d \geq 2$ and every sufficiently large $t$,
\[
\TR(\Gamma_{d,t}) > \R(\Gamma_{d,t}),
\]
and
\[
\liminf_{t \to \infty} \frac{\TR(\Gamma_{d,t})}{\R(\Gamma_{d,t})} \geq \frac{2d}{3}.
\]
Consequently,
\[
\sup \left\{ \frac{\TR(G)}{\R(G)} : G \text{ is connected and bipartite} \right\} = \infty.
\]
More precisely, for every $M > 0$ there exists a degree bound $\Delta_M$ and an infinite family of connected bipartite graphs of maximum degree at most $\Delta_M$ for which $\liminf \TR/\R \geq M$.
\end{theorem}

\begin{proof}
By Lemma~\ref{lem:count} and \eqref{eq:dgrid},
\[
\TR(\Gamma_{d,t}) \geq d\, t^{d-1}(t-1) + 1 = (d + o(1)) t^d.
\]
Combining this with \eqref{eq:ramseydgrid} gives the asserted ratio. Since $2d/3 > 1$ for every $d \geq 2$, the strict inequality follows for all sufficiently large $t$. Given $M > 0$, choose $d$ with $2d/3 \geq M$ and take $\Delta_M = 2d$.
\end{proof}


\begin{thebibliography}{99}
\bibitem[MSST]{MSST} G.~O.~Mota, G.~N.~S\'ark\"ozy, M.~Schacht, and A.~Taraz, \emph{Ramsey numbers for bipartite graphs with small bandwidth}, European Journal of Combinatorics \textbf{48} (2015), 165--176. \href{https://doi.org/10.1016/j.ejc.2015.02.018}{doi:10.1016/j.ejc.2015.02.018}.
\end{thebibliography}
\end{document}
